\section{The result and its context}

Foulkes' conjecture compares two ways of composing symmetric powers.
For a finite-dimensional complex vector space \(V\) and integers
\(1\le a\le b\), it asks for a \(\GL(V)\)-equivariant injection
\[
             \Sym^a(\Sym^b V)\hookrightarrow\Sym^b(\Sym^a V).
\]
Over \(\CC\), complete reducibility makes this equivalent to an
inequality between the multiplicities of every irreducible polynomial
representation. Write \(\Schur_\lambda V\) for the Schur module indexed
by a partition \(\lambda\), and \(s_\lambda\) for its character.
If \(h_r\) denotes the complete homogeneous symmetric function, then
\(h_a[h_b]\) is the stable character of \(\Sym^a(\Sym^b V)\), where
brackets denote plethysm. The conjecture says that
\(h_b[h_a]-h_a[h_b]\) is Schur-positive: all its Schur coefficients
are nonnegative integers. We prove the sixth case.

\begin{theorem}\label{thm:main}
For every integer \(b\ge6\) and every finite-dimensional complex vector
space \(V\), there is a \(\GL(V)\)-equivariant injection
\[
                 \Sym^6(\Sym^b V)\ \hookrightarrow\ \Sym^b(\Sym^6 V).
\]
Equivalently, \(h_b[h_6]-h_6[h_b]\) is Schur-positive for every \(b\ge6\).
\end{theorem}

Theorem~\ref{thm:main} resolves this case of Foulkes' conjecture
positively, with no restriction on \(\dim V\). The diagonal case
\(b=6\) is immediate; the content is the comparison for every larger
\(b\).

\paragraph{History and the canonical map.}
The conjecture originates in Foulkes' study of concomitants
\cite{Foulkes1950}. For binary forms, classical Hermite reciprocity
gives an isomorphism between the two symmetric powers for all \(a,b\);
see \cite[Introduction]{RaicuSamWeyman2022}. In higher dimension one
seeks a multiplicity inequality instead.
The case \(a=1\) is immediate. Thrall's decomposition formulas imply
\(a=2\) \cite{Thrall1942}; see also \cite[p.~237]{DentSiemons2000}.
Dent and Siemons proved \(a=3\) \cite{DentSiemons2000}.

A natural approach uses the canonical Foulkes--Howe map
\[
 \Psi_{a,b,V}:\Sym^a(\Sym^b V)\longrightarrow\Sym^b(\Sym^a V).
\]
Realize the source as symmetric tensors in \(V^{\otimes ab}\), arrange
the tensor positions in \(a\) rows of length \(b\), transpose the array,
and symmetrize the resulting blocks. McKay proved that injectivity
for every \(V\) at a seed \((a,b_0)\), with \(b_0\ge a\), propagates to every
\((a,b)\) with \(b\ge b_0\)
\cite{McKay2008,Ikenmeyer2015}.
The computation at \((4,4)\) by M\"uller and Neunh\"offer therefore
yields \(a=4\) \cite{MuellerNeunhoeffer2005};
Cheung, Ikenmeyer, and Mkrtchyan's injectivity computation at \((5,6)\)
yields \(a=5\), with the diagonal comparison again immediate
\cite[preprint, Theorem~6(a)]{CheungIkenmeyerMkrtchyan2017}.
Canonical injectivity is stronger than Foulkes' multiplicity comparison.
The canonical maps at \((5,5)\) and \((6,6)\) have nonzero kernels in
suitable dimensions, although their source and target representations
coincide \cite{MuellerNeunhoeffer2005,CheungIkenmeyerMkrtchyan2017}.
For the sixth diagonal this already occurs when \(\dim V\ge6\)
\cite[preprint, Theorem~6(c)]{CheungIkenmeyerMkrtchyan2017}.
Our finite comparison does not require injectivity of a prescribed map.

Exact character calculations give another approach.
Evseev, Paget, and Wildon verified Foulkes' multiplicity inequalities
when the two parameters sum to at most \(19\), including the sixth
case through \(b=13\)
\cite[Corollary~5.2]{EvseevPagetWildon2014}.
Their character-deflation recurrence is a predecessor of the exact
source calculation used here. The present certificates combine such
calculations with bounds on entire families of multiplicities, leaving
only a small band for exact comparison.

Brion proved eventual truth for fixed \(a\)
\cite[Corollary~1.3]{Brion1993} and later gave a dimension-dependent
bound for surjectivity of the canonical map in the opposite direction
\cite[Theorem~3.3]{Brion1997}.
That map, \(\Psi_{b,a,V}\), is realized as multiplication in the invariant
ring associated with the normalization of the Chow variety
\cite[Section~7]{Landsberg2015}; \cite[Section~2]{RaicuSamWeyman2022}.
The companion paper proves its surjectivity for \(b\ge a(a-1)\),
uniformly in dimension \cite[Theorem~1]{OpenAIQuadratic2026}.
At \(a=6\), an invariant complement to its kernel supplies an
embedding for \(b\ge30\). The remaining task is to compare all
multiplicities for \(6\le b\le29\).

For binary forms, recent generalized comparisons concern plethysms
with different parameter pairs: Raicu, Sam, Weyman, and Yang prove
maximal rank under divisibility conditions
\cite[Theorem~1.4]{RaicuSamWeymanYang2026}, and Gangl, Guti\'errez,
and Szwej identify the dual map through a substitution construction
\cite[Theorem~1.2]{GanglGutierrezSzwej2026}.
These results concern a two-dimensional underlying space; the theorem
above treats every finite dimension.

\paragraph{Proof strategy.}
The source \(\Sym^6(\Sym^b V)\) has Schur support of length at most
six. Let \(F(\lambda)\) and \(Q(\lambda)\) denote its multiplicity and
the corresponding target multiplicity, respectively, for a partition
\(\lambda\) of \(6b\). Section~\ref{sec:reduction} reduces the problem
to these six-part partitions and supplies rectangular highest-weight
shifts of known comparisons and a chart certificate
for partitions with \(\lambda_2+\cdots+\lambda_6\le b\).
The canonical multiplication map also proves the main large-degree
range in Corollary~\ref{cor:quadratic-range}.

The finite argument compares upper bounds for \(F\) with lower bounds
for \(Q\). Section~\ref{sec:upper} bounds source multiplicities by
counting possible intermediate partitions in a signed strip recurrence,
and gives exact recurrences for both characters.
Section~\ref{sec:lower} amplifies known target multiplicities:
a basis of highest-weight polynomials can be chosen with distinct leading
monomials, so a sumset bound supplies many independent products of those
polynomials.
Multiplication by a single nonzero highest-weight polynomial then
transfers these lower bounds to further weights.
Section~\ref{sec:certificates} organizes the resulting tests into a
finite computation. It checks every partition for \(6\le b\le25\).
In increasing degrees from \(26\) onward, rectangular shifts reduce
the remaining problem to \(\lambda_6<6\); the certificates leave only
partitions in the band \(\lambda_3+\cdots+\lambda_6\le27\).
Section~\ref{sec:band} computes exact differences throughout this band
by reusing seven truncated rational-function numerators.
Section~\ref{sec:verification} proves the arithmetic bounds and
assembles the range coverage.

An independent route is retained in full. The chart calculation used
in the finite comparison, combined with polarization, proves canonical
surjectivity for \(b\ge150\) in Proposition~\ref{prop:large}.
The polarization argument is related to the successive weight shifts
in the proof of McKay's theorem \cite[Section~5]{Ikenmeyer2015}.
Both certificate programs run through \(b=149\), so the extended finite
verification and this chart bound prove Theorem~\ref{thm:main} without
the quadratic-stabilization input. The main route needs only the
finite comparisons through \(b=29\).
The complete programs are printed in the appendices and supplied with
reference outputs and a reproduction script. The formulas explain
what the programs compute; their exact execution supplies the finite
nonnegativity assertions.

Throughout, we use characteristic-zero complete reducibility and
highest-weight theory for polynomial \(\GL\)-representations
\cite[Theorem~9.19 and Proposition~14.13]{FultonHarris1991},
Schur functors and the Cauchy decomposition
\cite[Theorem~6.3 and Exercise~6.11(b)]{FultonHarris1991},
and the Schur alternant formula \cite[Chapter~I, (3.1)]{Macdonald1995}.
Stable Schur coefficients are unchanged on specializing extra character
variables to zero whenever the partition length remains admissible
\cite[Chapter~I, (3.2)--(3.3)]{Macdonald1995}.
The character interpretation of plethysm is recalled in
\cite[Chapter~I, Appendix~A, Sections~7--8]{Macdonald1995}.
We also write \(e_r\) for the elementary symmetric function and set
\(h_0=e_0=1\).
